% !TEX encoding = UTF-8 Unicode
% !TEX TS-program = XeLaTex


\documentclass[8pt, twoside, a5paper]{extreport}
\usepackage{lipsum} % ----------------------------------------------------------- cancellabile
\usepackage{blindtext}	% dummy text

\usepackage{latexsym}
\usepackage[polutonikogreek, italian, english]{babel}
\usepackage[T1]{fontenc}
%\usepackage[utf8]{inputenc}

\usepackage[margin=1cm]{geometry}

\usepackage{multicol}
\usepackage[hang, small,labelfont=bf,up,textfont=it,up]{caption} % Custom captions under/above floats in tables or figures
\usepackage{booktabs} % Horizontal rules in tables
\usepackage{float} % they need to be placed in specific locations with the [H] (e.g. \begin{table}[H])

\usepackage{graphicx}
\usepackage{amssymb}
\usepackage{epstopdf}

\usepackage{hyperref} % For hyperlinks in the PDF

\usepackage{titlesec} % Allows customization of titles
\renewcommand\thesection{\Roman{section}} % Roman numerals for the sections
\renewcommand\thesubsection{\Roman{subsection}} % Roman numerals for subsections
\titleformat{\section}[block]{\large\scshape\centering}{\thesection.}{1em}{} % Change the look of the section titles
\titleformat{\subsection}[block]{\large}{\thesubsection.}{1em}{} % Change the look of the section titles

% \usepackage{fancyhdr} % Headers and footers
% \pagestyle{fancy} % All pages have headers and footers
% \fancyhead{} % Blank out the default header
% \fancyfoot{} % Blank out the default footer
% \fancyhead{}%[C]{Sei volte EMUfest $\bullet$ Novembre 2013 }%$\bullet$ Vol. XXI, No. 1} % Custom header text
% \fancyfoot[RO,LE]{\thepage} % Custom footer text

\renewcommand{\thefootnote}{\textasteriskcentered}



\usepackage{mwe,tikz}
\usepackage{tcolorbox}

\usepackage{color}

\usepackage{fontspec,xltxtra,xunicode}
\defaultfontfeatures{Mapping=tex-text}
\setromanfont[Mapping=tex-text]{Source Sans Pro}
\setsansfont[Scale=MatchLowercase,Mapping=tex-text]{Source Sans Pro}
\setmonofont[]{Source Code Pro}

\newfontfamily{\svolk}{Volkhov}

\linespread{1.2}

\usepackage{lipsum}

\definecolor{supercolor}{RGB}{3,39,36}

%\topmargin -2.66in

\newtcbox{\mybox}{nobeforeafter,colframe=supercolor,colback=supercolor,boxrule=0.5pt,arc=8pt,
  boxsep=0pt,left=2pt,right=2pt,top=6pt,bottom=2pt,tcbox raise base}

% !TEX encoding = UTF-8 Unicode
% !TEX TS-program = XeLaTex
% !TEX root = EMU2015_booklet.tex

%----------------------------------------------------------------------------------------
%	NEW COMMANDS 2015
%----------------------------------------------------------------------------------------

\newcommand{\greco}[1]{%
\begin{otherlanguage*}{greek}#1\end{otherlanguage*}}


%----------------------------------------------------------------------------------------
\newcommand{\livel}[6]{%
\noindent \textsc{#1}\\ %
\noindent \textbf{\textit{#2}} -- #3\\%
\noindent #4\\ %
\noindent #5 -- \textsc{#6}%
\\
}%

%: #1 autore, #2 titolo, #3 anno ed esecuzione, #4 tipologia, #5 strumenti, #6 esecutori

%\livel{}
%{}{}
%{}
%{}{}


%-----------------------------------------------------------------------------------------------------------------
\newcommand{\acusmatico}[4]{%
\noindent \textsc{#1}\\ %
\noindent \textbf{\textit{#2}} -- #3\\%
\noindent #4%
\\
}%

%: #1 autore, #2 titolo, #3 anno ed esecuzione, #4 tipologia
%\acusmatico{}
%{}{}
%{}

%-----------------------------------------------------------------------------------------------------------------
\newcommand{\acusmatici}[3]{%
\noindent \textsc{#1}\\ %
\noindent \textbf{\textit{#2}} -- #3%
\\
}%

%: #1 autore, #2 titolo, #3 anno ed esecuzione
%\acusmatici{}
%{}{}


%-----------------------------------------------------------------------------------------------------------------
\newcommand{\descrizione}[2]{%
\noindent \textbf{\textit{#1}} %
#2 %
\\
}%

%: #1 titolo, #2 abstract
% \descrizione{}{}


%-----------------------------------------------------------------------------------------------------------------
\newcommand{\biografia}[2]{%
\noindent \textsc{#1} %
#2 %
\medskip
}%

%: #1 autore, #2 biografia
% \biografia{}{}



%----------------------------------------------------------------------------------------
%	TITLE SECTION
%----------------------------------------------------------------------------------------

\title{
	\svolk{CONSERVATORIO DI MUSICA S. CECILIA} \\
	%\vspace{-15mm}
	\fontsize{50}{50}
	\svolk{
		\emph{
			EMUFest 2015
			}
		}
	} % Article title

\author{
	\textsc{Conservatorio di Musica Santa Cecilia} \\
	5 -- 10 ottobre 2015 \\
	\textsc{Università di Roma Tor Vergata} \\
	13 e 14 ottobre 2015 \\
	Roma
}

\vfill

\date{}

%\dedica{.2\textwidth}{\small A Monica:\\
%per tutto ciò che mi hai insegnato\\
%e per tutto ciò che ancora avresti dovuto insegnarmi.}

%----------------------------------------------------------------------------------------

\begin{document}
\pagestyle{empty}
\maketitle 


%----------------------------------------------------------------------------------------

\section*{5 ottobre 2015 -- ore 17:00}

\subsection*{{\small Concerto Acusmatico in Cupola Ambisonic} \\
	\textsf{Aula Bianchini}}

{\fontsize{30}{30} \svolk{\emph{ATTO I}}}

\subsection*{\textsf{regia del suono di Francesco Ziello}}

\bigskip

\begin{multicols}{2}

% !TEX encoding = UTF-8 Unicode
% !TEX TS-program = XeLaTex
% !TEX root = EMU2015_booklet.tex

\acusmatici{Ursula Meyer-K\"onig}
{Allears}{2012-13}{8'}

\acusmatici{Benjamin O'Brien}
{Along the eaves}{2012-13}{8'20''}

\acusmatici{Dennis Deovides A. Reyes III}
{Bolgia}{2014}{7'31''}

\acusmatici{Dimitrios Savva}
{Balloon Theories}{2012-13}{14'30''}

\acusmatici{Jones Margarucci}
{Inhabitated Places\textunderscore Part II}{2012-13}{5'52"}


%\descrizione{Allears}{The inspiration for this work originally came from a series of intensive discussions with people who are deaf or have hearing impairments. We talked about the pros and cons of technical apparatuses such as hearing aids or cochlea implants, the different ethical and emotional responses people have to them, and the identity issues they raise. Wearing hearing aids also changes how sounds are perceived, sometimes causing interference, distortions, diminished spatial perception and noise overflow.}

\descrizione{Allears}{Originariamente l'ispirazione per questo lavoro  proveniva da una serie di intense discussioni con persone non udenti o che hanno problemi d'udito. Abbiamo parlato dei pro e dei contro di apparati tecnici, come apparecchi acustici o impianti cocleari, le diverse risposte etiche ed emotive che le persone sentono, e i problemi d'identità che sollevano. 
Indossare apparecchi acustici cambia anche come i suoni vengono percepiti, a volte causando interferenze, distorsioni,percezione spaziale ridotta e troppo pieno di rumore.}


%\descrizione{Along the Eaves}{takes its name from the following line in Franz Kafka’s “A Crossbreed” “On the moonlight nights its favorite promenade is along the eaves.” To compose the work, I developed custom software and used these programs in different ways to process and sequence my source materials, which, in this case, included audio recordings of water, babies, and string instruments. My interest is to create sonic coincidences that suggest relationships between sounds and the illusions they foster.}

\descrizione{Along the Eaves}{ prende il nome dalla riga che di Franz Kafka "Incrocio". "Sulle notti di luna la sua passeggiata preferita è lungo la grondaia" Per comporre l'opera, ho sviluppato software personalizzato e utilizzato questi programmi in modi diversi per elaborare e sequenziare i miei materiali di base, che, in questo caso, include registrazioni audio d’acqua, bambini, e di strumenti a corda. Il mio interesse è quello di creare coincidenze sonore che suggeriscono i rapporti tra i suoni e le illusioni che promuovono.}


%\descrizione{Bolgia}{Bolgia is an Italian word that means pocket or trench.  This term has been used by Dante Alighieri in his notable literary work Inferno.  Bolgia is a stereo fixed media electroacoustic composition, which depicts Alighieri’s journey to the eighth circle of hell, and his experiences to its horrific environment.  The musical gestures and sonic events of the piece evoke the different sounds and emotions of hell.}

\descrizione{Bolgia}{è una parola italiana che significa "fossa" e "luogo chiassoso in cui regna la confusione". Questo termine è stato usato da Dante Alighieri nel suo noto lavoro letterario "Inferno". Bolgia è un brano stereofonico fisso per composizioni elettroacustiche, che illustra il viaggio di Alighieri nell'ottavo girone dell'inferno, e la sua esperienza in questo posto terribile. I gesti musicali e l'evento sonoro del pezzo evocano i diversi suoni e le diverse emozioni dell'inferno.}


%\descrizione{Balloon Theories}{«I was always enjoying squeezing balloons, pressing them with my fingers until they pop… It has not been up until now that I realized why…»}

\descrizione{Balloon Theories}{«Ho sempre trovato divertente strizzare palloncini, premerli con le dita fino allo scoppio... Non mi è mai interessato fino a quando non ho capito perché...»}


%\descrizione{Inhabitated Places\textunderscore Part II}{Inhabitated Places part II is based on the concept of algorithmic composition. Although the general shape of this piece has been determined in a conventional way, every sound that one can hear are selected in real time by different algorithms written in SuperCollider. These algorithms choose randomly audio files from different folders and play them at different speeds and in different moments. It is as if we had placed several different objects in several boxes (that represent our shape), but every time we open one of these boxes the objects placed inside are positioned differently from how we had left them previously. The piece was also rendered in B-format Ambisonic.}

\descrizione{Inhabitated Places\textunderscore part II}{ è una composizione elettroacustica basata sul concetto di musica algoritmica. Sebbene la forma generale del brano sia stata determinata apriori e in modo convenzionale, tutti i suoni che ascoltiamo vengono scelti in tempo reale da vari algoritmi scritti in SuperCollider. Questi algoritmi selezionano in modo pseudocasuale dei samples da diverse cartelle e li riproducono a velocità diverse e in diversi momenti.  
È come se avessimo sistemato in una scatola (che in questo caso rappresenta la struttura dell’opera) degli oggetti in un dato ordine, ma ogni qual volta apriamo la scatola li troviamo disposti in modo differente da come li avevamo lasciati.}


\end{multicols}

\clearpage
%----------------------------------------------------------------------------------------

\section*{5 ottobre 2015 -- ore 20:30}

\subsection*{{\small Concerto Acusmatico \& Live Electronics\footnote{ Il concerto verrà trasmesso in diretta streaming da Radio Cemat}} \\
	\textsf{Sala Accademica}}

{\fontsize{30}{30} \svolk{\emph{ATTO II}}}

\subsection*{\textsf{Regia del suono di Federico Paganelli}}

\bigskip

\begin{multicols}{2}

%% !TEX encoding = UTF-8 Unicode
% !TEX TS-program = XeLaTex
% !TEX root = EMU2015_booklet.tex

\livel{Simone Cardini}
{Potlach}{2014 - vincitore Premio Bucchi} %9'30''
{per flauto basso ed elettronica su supporto}
{flauto basso}{Gianni Trovalusci}

\livel{Dominique Schafer}
{Cendre}{2008} %10'00''
{per flauto basso e live electronics}
{flauto basso}{Gianni Trovalusci}

\livel{Maria Cristina De Amicis}
{[re-spì-ro]}{2015 - prima esecuzione assoluta} %8'00''
{per flauti e supporto digitale}
{flauti}{Gianni Trovalusci}

\acusmatico{Christian Eloy}
{La cicatrice d'Ulysse}{new version 2015} %13'00''
{acusmatico}

\livel{Giorgio Nottoli}
{7 Isole}{2015 - prima esecuzione assoluta}%15'
{per flauto, percussioni e live electronics}
{flauto}{Gianni Trovalusci}
percussioni -- \textsc{Gianluca Ruggeri}
\\



%\descrizione{Potlach}{was a ritual ceremony performed among certain Native American tribes, during which the members of the community exchanged and squandered gifts, contributing to the cohesiveness of the social concord, whithin a totally inverted logic with respect to our free economy, following a reciprocity principle. The formal juxtaposition of the piece, disperses at once those elements that will constitute later on the concealed foundation that penetratrates the work.}

\descrizione{Potlach}{Il Potlach era una cerimonia rituale d’alcuni popoli di Nativi Americani, durante la quale i membri della tribù scambiavano e dilapidavano doni, in una logica completamente invertita rispetto alla nostra economia. La giustapposizione del brano \textit{dilapida} da subito quegli elementi che costituiranno il substrato che pervade l'opera. È la simbolizzazione dell'immagine che ne permette un'esegesi palesata; l'interprete realizza sé attraverso la sublimazione di gesti e momenti performativi, rappresentando la fenomenologia degli stessi. Da qui la necessità di superare il momento poietico ed estesico \textit{donando} la performance per un recupero del sistema di relazioni come tale.}

%\descrizione{Cendre}{draws on the idea from the impermanence of materiality. The title of the piece gives further meaning referring to a fragile and delicate state, but also the potential of ashes as a fertilizer for something new. The bass flute being suspended within the electronics extends space and timbre, at times detaching itself within its own identity, but repeatedly gets reabsorbed into space. The piece was written for and premiered by Mario Caroli during his Fromm Residency at Harvard University.}

\descrizione{Cendre}{prende ispirazione dall’idea della precarietà di ciò che è materiale. Il titolo del brano acquista ulteriore significato riferendosi ad uno stato di fragilità e delicatezza, ma anche alle potenzialità fertilizzanti della cenere per qualcosa di nuovo. Grazie alla sospensione del flauto basso attraverso l’elettronica, il suo spazio ed il suo timbro vengono estesi, a tratti allontanandosi dalla propria identità per poi essere riassorbito all’interno dello spazio. Il brano è stato composto ed eseguito in prima assoluta da Mario Caroli durante la Fromm Residency presso la Harvard University.}

%\descrizione{[re-spì-ro]}{Re-spi-ro (Breath) in this work, is the alternation of air movements which could produce real rhythmic and tonal structures. The sound of each breath adds and interacts itself with the previous one, drawing a music form which contracts and expands itself as well. Through the rhythmic articulation, inhalation and exhalation, the breath transforms and builds itself, removes and contradicts what has been just expressed. The digital support creates a constructive dialectic, supporting gently the concrete material. Breath transformation is described by the conduct of three different changing elements: localization, movement (direction and speed) and space size. In this way , the auditive perception of the space is transposed to the visual perception level. In this work I tried to realize the interaction between these changing elements on the acoustic level and associating with tones and pitches the breath expressive gestures of the performer.}

\descrizione{[re-spì-ro]}{Il re-spi-ro, in questa opera, è inteso come l’alternanza dei movimenti dell’aria in grado di generare vere e proprie strutture ritmiche e timbriche. Il suono di ogni respiro si aggiunge e interagisce con i precedenti, disegnando una forma musicale che si contrae e si espande.  Attraverso l’articolazione ritmica, l’inspirazione e l’espirazione, il respiro si trasforma e costruisce, dissolve o contraddice ciò che è stato appena espresso. Il supporto digitale crea una dialettica costruttiva sostenendo delicatamente il materiale concreto. La trasformazione del respiro è descritta attraverso il comportamento di tre variabili: localizzazione, movimento (direzione, velocità) e dimensione dell’ambiente. In questo modo la percezione uditiva dello spazio viene trasposta nel dominio della percezione visiva. In questo brano ho cercato di realizzare l’interazione tra queste variabili nel dominio acustico associando con timbri e altezze i gesti espressivi del respiro dell’esecutore.}

%\descrizione{La cicatrice d'Ulysse}{This title, borrowed from Erich Auerbach, German writer and critic deceased in 1957, immediately sets the scene on a plane where realism is depicted in the aesthetics of Western music. Electroacoustic music has the ability, using “concrete” sounds (with all the ambiguity implied by the word), of giving us an immediate sensation of reality, which we can all situate in relation to our references and private objects of reference; these are the profound notions of sublimitas and humilitas, which merge and unite in musical expression. Every listener will be able to decipher his or her own images of a collective mental universe from the essence of this kind of artistic creation, just like the scar by which Ulysses was recognised by his former servant woman.}

\descrizione{La cicatrice d'Ulysse}{Il titolo, preso in prestito da Erich Auerbach, scrittore e critico tedesco morto nel 1957, ambienta immediatamente la scena all’interno di un aereo, tratteggiandola con un realismo tipico dell’estetica musicale occidentale. Attraverso l’uso di suoni concreti (con tutte le ambiguità insite nella parola stessa), la musica elettroacustica possiede l’abilità di restituirci una percezione immediata della realtà, che è possibile collocare in relazione ai propri riferimenti e agli oggetti segreti cui intende riferirsi; questi sono i significati profondi di sublimitas e humilitas, che si fondono e riuniscono nell’espressione musicale. Dall’essenza di questo tipo di creazione artistica ogni ascoltatore è in grado di decodificare la propria immagine di un universo mentale collettivo, proprio come la cicatrice che permise alla serva di riconoscere Ulisse.}

%\descrizione{7 Isole}{I think the idea of the island is a fascinating metaphor, because I realize to think about many important things  as they are formed, in fact, by islands, which each contain a different world and a different feel, but together form a unitary context. In "7 Isole", each "island" is characterized by a particular combination of movement, pitches and colors of the sound. It consists of seven small pieces separated from one another, that can be performed in any order, however, strongly linked by way of forming the sound material and by a unitary constructive thinking. The instruments are used both with extended techniques that traditional, allowing different nuances both for the continuous sound as well as impulsive, both for what determined pitch and for what similar to the noise. The electronics supports and extends the sound of the instruments and, where necessary, becomes the instrument itself, completing the construction of the sound field. The dynamic distribution  in the listening space is obtained according to the principles of the method Ambisonic. Some of the processing and sound synthesis used have been developed by the author.}

\descrizione{7 Isole}{Quella dell'isola è per me una metafora affascinante, in quanto mi accorgo di pensare a molte cose importanti come fossero costituite, appunto, da isole, che contengono ciascuna un mondo diverso e quindi un diverso sentire, ma, insieme, costituiscono un contesto unitario. In 7Isole, ciascuna "isola" è caratterizzata da una particolare combinazione di movimento, altezze e colori del suono. Si tratta di sette piccoli pezzi fra loro separati, che si possono eseguire in un qualsiasi ordine, tuttavia fortemente legati dal modo di formare il materiale sonoro e da un pensiero costruttivo unitario. Gli strumenti sono utilizzati sia con tecniche estese che tradizionali, consentendo diverse sfumature sia per il suono continuo che per quello impulsivo, sia per quello ad altezza determinata e per quello simile al rumore. L'elettronica sostiene ed  estende il suono degli strumenti e, dove necessario, diviene strumento essa stessa, completando la costruzione del campo sonoro. La distribuzione dinamica nello spazio d'ascolto è ottenuta secondo i principi  del metodo Ambisonic. Alcune delle elaborazioni e sintesi del suono utilizzate sono state sviluppate dall'autore.}











\end{multicols}

\clearpage

%----------------------------------------------------------------------------------------

\section*{6 ottobre 2015 -- ore 17:00}

\subsection*{{\small Concerto Acusmatico in High Order Ambisonic Domain} \\
	\textsf{Aula Bianchini}}

{\fontsize{30}{30} \svolk{\emph{ATTO III}}}

\subsection*{\textsf{regia del suono di Federico Paganelli}}

\bigskip

\begin{multicols}{2}

%% !TEX encoding = UTF-8 Unicode
% !TEX TS-program = XeLaTex
% !TEX root = EMU2015_booklet.tex

\acusmatici{Jan Jacob Hofmann}
{Horizontal and Vertical Lines}{2005}{7'53''}

\acusmatici{Clemens Von Reusner}
{Topos Concrete}{2014}{9'18''}

\acusmatici{Damián Anache}
{Capturas del Único Camino}{2014-15}{14'29''}

\acusmatici{Jan Jacob Hofmann}
{Tensile Elements}{2000}{9'30''}

% \descrizione{Horizontal and Vertical Lines}{Two kinds of elements mark the context of the piece: Sounds of extremely long duration in the distance, shimmering and sharp like steel, stretched out in the timeline. Generated by patterns of chaotic oscillation and being non-linear, these horizontal elements refer to infinity. Their vertical organisation, the proportion of pitch among each other is organised by the harmonic ratio though. Much closer, several impulses: Sounds of infinite shortness, containing the whole spectra of frequencies in a mathematical sense, but having no extension in time. Being strictly vertical, they contradict with the long stretched horizontal elements. These two sonic elements, considered as horizontal and vertical, set the background structure of which matter is generated. The sound for this composition is derived from a non-linear algorithm for sound generation by physical modelling, as from granular synthesis. All the sounds have been created using Csound. The work was commissioned by the CRM for Musica Scienza 2005 where it got executed in June 2005.}

\descrizione{Horizontal and Vertical Lines}{Due elementi determinano la base del pezzo: suoni di eccezionale lunghezza dalla distanza, scintillante e tagliente come l'acciaio. Si estendevano lungo l'asse del tempo. Generati dal modello di oscillazione caotica e non-linearità, questi elementi orizzontali mostrano all'infinito. La loro organizzazione verticale, cioè il loro rapporto fra loro è pero disposte secondo le leggi dell'armonia. Molto più vicino: differenti suoni che si assomigliano le impulsi: suoni di durata infinitamente breve, tuttavia, contengono l'intero spettro di frequenze in senso matematico, ma non hanno estensione nel tempo. Strettamente allineati verticalmente, in contrasto con la lunga tesa elementi orizzontali. Entrambi questi elementi sonori, la forma orizzontale e verticale, creano la struttura basale di quella si sviluppano le forme. I suoni di questa composizione sono generati da un algoritmo non-lineare, così come la sintesi granulare. I suoni sono stati generati con il programma Csound. Il brano "Linee Orizzontale e Verticale" è stata composto come contributo a "Musica e Scienza" nel 2005, per conto di CRM, Roma.}

% \descrizione{Topos Concrete}{Concrete is a building material, a kind of unshaped dry powder made of sand, granulated stones and cement, dusty and chaotic. Mixed with water it becomes flexible and fluid and goes into a metamorphosis to become dry again, static and resistable and of any wanted shape. Aspects of working with native granularity, fluidness as well as stiffness and different kind of acoustic spaces were leading ideas of the composition.}

\descrizione{Topos Concrete}{Il calcestruzzo ("Concrete") è un materiale da costruzione, una sorta di polvere secca informe costituita di sabbia, pietre e cemento granulati, polveroso e caotico. Impastato con l'acqua diventa flessibile e fluido, quindi torna ad essere di nuovo asciutto, statico e resistente e di qualsivoglia forma. La composizione del brano è guidata da alcuni aspetti quali il lavorare con la granularità e la fluidità così come con la rigidità e con diversi tipi di spazi acustici.}

% \descrizione{Capturas del Único Camino}{"First Landscape" is a movement of “Capturas del Único Camino”, a piece which involves generative art ideas for offering an attractive object of passive contemplation. It is Ambisonics encoded and made with acoustic instruments samples recorded and performed by the composer. Then the samples are handled by a Pure Data patch according to the random events score of the piece. So that, the algorithm works as an electronic performer of the piece. info: http://conceptocero.com/capturasdelunicocamino}

\descrizione{Capturas del Único Camino}{"First Landscape" è un movimento di "Capturas del Único Camino", un brano che impiega le idee dell'arte generativa, per offrire un interessante oggetto di contemplazione passiva. Il brano è codificato Ambisonics e realizzato con campioni di strumenti acustici registrati e suonati dal compositore. I campioni sono gestiti da una patch di Pure Data, seguendo una partitura di eventi randomici. In questo modo l'algoritmo funziona come se fosse l'esecutore elettronico del pezzo. info: http://conceptocero.com/capturasdelunicocamino}

% \descrizione{Tensile Elements}{Subject of \textit{Tensile Elements} are special kinds of elements, characterised by their quality of sound, their movement and interaction. Elements pass through the space, defining it, having almost materiality but still being bodiless. Concentrations and disintegrations happen, the elements seem to have almost their own way of behaviour.}

\descrizione{Tensile Elements}{Tema del soggetto sono elementi che si caratterizzano per il loro determinato suono teso, così come per i loro movimenti e le loro interazioni. Tali elementi si muovono attraverso lo spazio, lo definiscono ed hanno quasi una materialità, nonostante ciò rimangono senza corpo. Addensamenti e dissolvenze si succedono, gli elementi sembrano avere una vita propria.}


%Sonic Architecture 

%Jan Jacob Hofmann's work is spatial electronic music, or, as he calls it,  “sonic architecture” – architecture made by sounds. Using the Ambisonic-Method makes it possible to assign a precise Position  in 3D space to every sound.  The sonic material of the composition generates a sonic space surrounding the listener. Sounding virtual materials do generate complex architectures made of sound, spaces in transformation do emerge.
%
%This is made possible by using a combination of the Ambisonic Method in combination with Csound, a computer-language for sound synthesis and digital signal processing. Jan Jacob Hofmann conceived and created a set of instruments that allow to assign a position in 3d space to every sound. Also the sound is placed in a sonic environment allowing distance perception by the creation of spatial reverb and early reflections, also using the ambisonic method.
%
%Embedded in an environment for composition, this code makes it possible to compose works of spatial music. This code is available open source on his site www.sonicarchitecture.de and is already used by other composers.
%
%Latest development is the extension of the spatial assignment of sounds to granular clouds so that spatial granular synthesis became possible allowing swarms and clouds of sound in space.










\end{multicols}

\clearpage
%----------------------------------------------------------------------------------------

\section*{6 ottobre 2015 -- ore 20:30}

\subsection*{{\small Concerto Acusmatico \& Live Electronics\footnote{ Il concerto verrà trasmesso in diretta streaming da Radio Cemat}} \\
	\textsf{Sala Accademica}}

{\fontsize{30}{30} \svolk{\emph{ATTO IV}}}

\subsection*{\textsf{Regia del suono di Luana Lunetta}}

\bigskip

\begin{multicols}{2}

%% !TEX encoding = UTF-8 Unicode
% !TEX TS-program = XeLaTex
% !TEX root = EMU2015_booklet.tex

\livel{Alessia Forganni, Daniele Vantaggio}
{KOMA - Kontrolled Organic Movement Act}{10'}
{voce, pianoforte,  kinekt, e live electronics}
{2015}

\brano{Massimo Massimi}
{Kronos in Fabula I}{10'}
{acusmatico}
{2015}\\

\livel{Claudia Jane Scroccaro, Federico Ripanti}
{On-y r\^eve}{7'40''}
{per pianoforte e nastro}
{2015}

\livel{Carlos D. Perales}
{17 haiku}{10'}
{flauto e live electronics}
{2012}

\livel{Massimiliano Mascaro, Giuseppe Zampetti}
{L'abile medico}{10'}
{per soprano, mezzosoprano e live electronics}
{2015}


\esecutore{voce, pianoforte}{Alessia Forganni}
\esecutore{pianoforte}{Daniele Buccio}
\esecutore{flauto}{Alessandro Pace}
\esecutore{soprano}{Elisabetta Braga}
\esecutore{mezzosoprano}{Virginia Guidi}

%\vspace{-5mm}

% \descrizione{KOMA - Kontrolled Organic Movement Act}{Our entire life is made by acts. Life is articulated day by day through repetition. Millions of automatisms assure us a cognitive and energetic saving. So life seems to be an automated system that is difficult to avoid or control, without a real consciousness of our actions. Maybe we should see our actions as a ritual gesture, instead of a simple repetition of the same act. Sometimes we should give time, meaning and attention to our acts to really control them. The performance includes tape, live piano, voices, live electronics and sound controlling by body movement through Kinect and Leap Motion system.}%abbiamo foto

\descrizione{KOMA - Kontrolled Organic Movement Act}{Tutta la nostra vita è fatta di gesti: grazie alla ripetizione di milioni di essi la vita è scandita ogni giorno. Forme di automatismi permettono un risparmio cognitivo ed energetico e alla maggior parte non si fa caso. La vita appare così un sistema automatizzato dal quale è difficile sottrarsi e avere reale controllo, uno stato di assenza di reale coscienza e consapevolezza. Per rendere i nostri gesti un atto consapevole bisogna vederli nella loro “ritualità” piuttosto che “ripetitività”, conferendovi significato e tempo,  focalizzandovi al massimo la nostra attenzione. La performance include nastro, piano e voci dal vivo, live electronics e controllo del suono tramite movimento con i sistemi di sensoristica Kinect e Leap Motion.}%abbiamo foto

% \descrizione{Kronos in Fabula I}{The acousmatic piece “Kronos in Fabula 1” is part of a series of compositions dedicated to the speculation about time. The material’s acoustic density represents the time trend; the possibilities of the listening points’ collocation on the time axis are connected to a feasible fruition “directionality” linked to the sequence and reiteration of sounds events.}

\descrizione{Kronos in Fabula I}{“Kronos in Fabula I” è l’acusmatico di un ciclo di composizioni dedicate alla speculazione sul tempo. La densità acustica del materiale rappresenta l’andamento del movimento temporale, le possibilità di collocazione di punti di ascolto sull’asse del tempo vengono associate a una possibile “direzionalità” dell’ascolto legata al concatenarsi e al reiterarsi degli eventi sonori.}

% \descrizione{On-y r\^eve}{The work is a “nocturnal representation” for piano and tape, suffused with suspended atmospheres and textures typical of oneiric visions. The piano part is mostly rhapsodic, while the acousmatic elements enhance and filter the resonances produced through the tonal pedal, conveying granulated processes of prepared piano. A contrast between darkness and light alludes to a chiaroscuro counterpoint, where each tile is devised manipulating and recomposing elements derived from Alexander Scriabin’s \textit{Poème-Nocturne} op. 61, in a tribute to commemorate the centenary celebration of his passing.}

\descrizione{On-y r\^eve}{The work is a “nocturnal representation” for piano and tape, suffused with suspended atmospheres and textures typical of oneiric visions. The piano part is mostly rhapsodic, while the acousmatic elements enhance and filter the resonances produced through the tonal pedal, conveying granulated processes of prepared piano. A contrast between darkness and light alludes to a chiaroscuro counterpoint, where each tile is devised manipulating and recomposing elements derived from Alexander Scriabin’s \textit{Poème-Nocturne} op. 61, in a tribute to commemorate the centenary celebration of his passing.}

% \descrizione{17 haiku}{Heirs of haikai, comic-themed and outgoing figuralisms in Japanese poetry, haiku was taking its own identity in the complex art of expression of consciousness. In them, senses are suggested, not the meanings. Haiku does not mean, because it says nothing. Indeed, haiku does not have an informative function. But, perhaps, to the expressive function, a suggestive function should be added.} %aggiungi pdf

\descrizione{17 Haiku} {Eredi degli haikai, fumetti a tema e figuralismi uscenti della poesia giapponese, gli haiku si sono guadagnati una propria identità nella complessa arte di espressione della coscienza. In essi si suggeriscono i sensi, non i significati. Haiku non significa, perché non dice nulla. Infatti, gli haiku non hanno una funzione informativa. Ma, forse, alla funzione espressiva bisogna aggiungere una funzione suggestiva.} %aggiungi pdf

% \descrizione{L'abile medico}{The piece is a prayer. Trying to go further, beyond the suffering it is a difficult task; and sometimes it causes reconsideration, but eventually the courage leads to victory. Il Sutra del Loto is one of the most important texts of the literature of Mahayana Buddhism. In the first section the breath, the primary form of meditation predisposes to concentration becoming  a prayer, the supplication of  the Mystic Law. Later, through meditation we realize the complexity of existence: an illusory chaos. Finally the resolute soul emerges, confirming definitively the depth of his Essence. To increase the listening ability and participation we recommend to close the eyes.}%aggiungere pdf

\descrizione{L'abile medico}{Il brano è una preghiera. Cercare di andare oltre, al di là della sofferenza è un compito difficile, e talvolta porta dei ripensamenti, ma alla fine il coraggio conduce alla vittoria. Il Sutra del Loto è uno dei testi più importanti della letteratura del Buddismo \textit{Mah y na}. Nella prima sezione il respiro, forma primaria di raccoglimento, predispone alla concentrazione trasformandosi in preghiera, invocazione della Legge Mistica. In seguito, attraverso la meditazione si realizza la complessità dell'esistenza: è il caos apparente. Finalmente l'anima risoluta riemerge, affermando definitivamente la profondità della sua Essenza. Per aumentare la capacità di ascolto e di partecipazione si consiglia di seguire ad occhi chiusi.}%aggiungere pdf




\end{multicols}

\clearpage

%----------------------------------------------------------------------------------------

\section*{7 ottobre 2015 -- ore 12:30}

\subsection*{{\small Concerto con Spazializzazione Interattiva in HOA} \\
	\textsf{Aula Bianchini}}

{\fontsize{30}{30} \svolk{\emph{ATTO V}}}

\subsection*{\textsf{regia del suono di Balandino Di Donato}}

\bigskip

\begin{multicols}{2}

%% !TEX encoding = UTF-8 Unicode
% !TEX TS-program = XeLaTex
% !TEX root = EMU2015_booklet.tex

\livel{Luciano Berio}
{Sequenza III}{9'} 
{per voce femminile spazializzata}
{1965}

\livel{Cathy Berberian}
{Stripsody}{11'30''} %
{per voce sola spazializzata}
{1966}

\descrizione{Sequenza III}{La voce porta sempre con sé un eccesso di connotazioni. Dal rumore più insolente al canto più squisito, la voce significa sempre qualcosa, rimanda sempre ad altro da sé e crea una gamma molto vasta di associazioni. In Sequenza III ho cercato di assimilare musicalmente molti aspetti della vocalità quotidiana, anche quelli triviali, senza però per questo rinunciare ad alcuni aspetti intermedi ed al canto vero e proprio. Per controllare un insieme così vasto di comportamenti vocali era necessario frantumare il testo e in apparenza devastarlo, per poterne recuperare i frammenti su diversi piani espressivi e ricomporli in unità non più discorsive ma musicali. Era cioè necessario rendere il testo omogeneo e disponibile al progetto che consiste, nelle sue linee essenziali, nell’esorcizzare l’eccesso di connotazioni componendole in un’unità musicale. Ecco il breve testo «modulare» di Markus Kutter per Sequenza III: \\ «Give me	a few words for a woman \\ to sing a truth allowing us \\ to build a house without worrying before night comes» \\ In Sequenza III l’enfasi è posta sul simbolismo sonoro di gesti vocali e talvolta visivi, sulle “ombre di significato” che li accompagnano, sulle associazioni e sui conflitti che essi suggeriscono. Per questa ragione Sequenza III può anche essere considerata come un saggio di drammaturgia musicale la cui storia, in un certo senso, è il rapporto fra l’interprete e la sua stessa voce. Sequenza III è stata scritta nel 1965 per Cathy Berberian. \\ Luciano Berio}

% \descrizione{Sequenza III}{(author's note) The voice carries always an excess of connotations, whatever it is doing. From the grossest of noises to the most delicate of singing, the voice always means something, always refers beyond itself and creates a huge range of associations. In Sequenza III I tried to assimilate many aspects of everyday vocal life, including trivial ones, without losing intermediate levels or indeed normal singing. In order to control such a wide range of vocal behaviour, I felt I had to break up the text in an apparently devastating way, so as to be able to recuperate fragments from it on different expressive planes, and to reshape them into units that were not discursive but musical. The text had to be homogeneous, in order to lend itself to a project that consisted essentially of exorcising the excessive connotations and composing them into musical units. This is the “modular” text written by Markus Kutter for Sequenza III. \\ Give me	a few words	for a woman \\ to sing	a truth allowing us \\ to build a house	without worrying \\ before night comes \\ In Sequenza III the emphasis is given to the sound symbolism of vocal and sometimes visual gestures, with their accompanying “shadows of meaning”, and the associations and conflicts suggested by them. For this reason Sequenza III can also be considered as a dramatic essay whose story, so to speak, is the relationship between the soloist and her own voice. Sequenza III was written in 1965 for Cathy Berberian. Luciano Berio}

\descrizione{Stripsody}{«Cathy Berberian era alla ricerca di un testo per una delle sue performance musicali e contemporaneamente stava sviluppando una sorta di mondo sonoro usando solo le onomatopee codificate dalla lingua dei fumetti. Gradualmente si convinse che queste azioni musicali non avevano bisogno della musica scritta; così, mentre Cathy incominciava a \emph{cantare} questi suoni, Carmi procedeva a \emph{scrivere} la partitura. I due aspetti del lavoro sono nati insieme, e la voce di Cathy ha dato più di un suggerimento grafico mentre l'impaginazione di Carmi ha fornito più di una soluzione vocale». \\ Umberto Eco}

%\descrizione{Stripsody}{“Berberian was looking for a text for one of her musical performances, and thought of developing a kind of sound world using only the onomatopoeic inventions of the comic strips. Gradually the idea grew that this musical action had no need of music; thus, while Cathy began to sing these sounds, Carmi went on to write the score. The two aspects of the work were born together, and Cathy’s voice contributed more than one graphic suggestion while Carmi’s imagination produced more than one vocal solution.” E. Carmi, Stripsody – Interpretazione vocale di Cathy Berberian. Testo introduttivo di Umberto Eco, Arco d’Alibert – Kiko Galleries, Roma – Houston (Texas) 1966.}


\end{multicols}

\clearpage
%----------------------------------------------------------------------------------------

\section*{7 ottobre 2015 -- ore 20:30}

\subsection*{{\small Concerto Acusmatico \& Live Electronics\footnote{ Il concerto verrà trasmesso in diretta streaming da Radio Cemat}} \\
	\textsf{Sala Accademica}}

{\fontsize{30}{30} \svolk{\emph{ATTO VI}}}

\subsection*{\textsf{Regia del suono di Marco De Martino}}

\bigskip

\begin{multicols}{2}

%% !TEX encoding = UTF-8 Unicode
% !TEX TS-program = XeLaTex
% !TEX root = EMU2015_booklet.tex

\livel{Elisabetta Capurso}
{Sezioni per organo spazializzato}{9'}
{per organo e live electronics}
{2015}

\livel{Cristiana Colaneri, Anna Terzaroli}
{Fantasia per due}{8'} 
{per sassofoni, percussioni ed elettronica}
{2015}

\brano{Demian Rudel Rey}
{Cenizas del Tiempo}{7'04''}
{acusmatico}{2015}\\

\livel{Gy\"orgy Ligeti}
{Volumina}{9'30''}
{per organo}
{1961-62, rivista 1966}

\livel{Alessandro Ratoci}
{Studio Mannaro}{11'30''}
{per sassofono baritono e live electronics}
{2013}

\esecutore{organo}{Antonella Barbarossa}
\esecutore{sassofoni}{Enzo Filippetti}
\esecutore{percussioni}{Valerio Cosmai}

% \descrizione{Sezioni per organo spazializzato}{Had many different changes in its general structure, the most important certainly the recent shift of the registers of the instrument. In the performance of the EMUFest 2015 the score has a further change for the presence of some electronic processing. For reasons of expression have been used some spatial sound systems, that fit perfectly to the sound of the acoustic instrument. Some microphones, amplifiers, some spatial sound's elaboration are amply present in the electronic score of the composition Sezioni.}

\descrizione{Sezioni per organo spazializzato}{Sezioni è opera che appartiene al periodo della prima attività compositiva. Dopo alcuni anni ha avuto alcuni cambiamenti nella struttura generale, un radicale cambiamento nella registrazione organistica.Il principio del mutamento di vista, simile a prisma ruotante è così la legenda della scrittura compositiva personale. La struttura generale ha tre sezioni, all'interno delle quali tre micro-sezioni hanno vita di improvvisazione. Nell'esecuzione di EMUFest2015 Sezioni ha un'ulteriore mutamento per la presenza di interventi elettronici. Per una  necessità di esigenze espressive sono stati  utilizzati alcuni sistemi di spazializzazione del suono  che si integrano perfettamente  con il suono originario dello strumento acustico. Microfonaggio, amplificazione, spazializzazione del suono sono gli interventi elettroacustici presenti nel processo di elaborazione sonora di Sezioni, suggeriti fondalmentalmente dalla raffinata registrazione organistica.}

% \descrizione{Fantasia per due}{“Fantasia per due” is among the EMUfest 2015 commissions destined to a couple of students of the Santa Cecilia Conservatory, one from the School of Composition, the other from the School of Electronic Music. The electronics uses the sound material produced by acoustic instruments, processing it and storing it on tape and it is built keeping in mind an equal and complementary relationship with the acoustic component. The latter is built starting from the indeterminate vibrations of a gong and from a single sound of the sax which, through gradual micro-intervals, conquers a wider frequency range interacting with percussions up to its full expansion. The conclusion is related to incipit.}

\descrizione{Fantasia per due}{Il brano è una commissione EMUfest, destinata ad una coppia di autori entrambi studenti del Conservatorio Santa Cecilia, uno proveniente dalla Scuola di Composizione, l'altro dalla Scuola di Musica Elettronica. L'elettronica utilizza il materiale sonoro prodotto dagli strumenti acustici, elaborandolo e fissandolo su tape ed è costruita nell'ottica di un rapporto paritetico e complementare con la composizione acustica. Questa è costruita partendo dalle vibrazioni indeterminate del gong e da un unico suono del sax, che gradualmente, anche con microintervalli, conquista un campo frequenziale più vasto, nell' interazione con le percussioni, fino alla libera espansione; la conclusione si ricollega all' incipit.}

% \descrizione{Cenizas del Tiempo}{Cenizas del Tiempo is an electroacoustic quadraphonic work, inspired by the idea that time ceases his ashes in our lives, gradually our being is consumed and the same thing happens with the materials and the sound objects. Also, expresses the experience of time in a city where everything occurs very quickly. In the work referential sounds of urban environments are perceived, moreover, samples from different ashtrays developed as more abstract sounds through processed and over-processed.}

\descrizione{Cenizas del Tiempo}{è un'opera quadrifonica elettroacustica, ispirata al tempo che porta le sue ceneri nella nostra vita, a poco a poco il nostro essere è consumato e la stessa cosa accade con i materiali e gli oggetti sonori. Inoltre, esprime l'esperienza del tempo in una città dove avviene tutto molto velocemente. In questo lavoro sono percepibili suoni riferiti ad ambienti urbani, e campioni di diversi posacenere che diventano suoni sempre più astratti attraverso continui processi.}

% \descrizione{Volumina}{eng}

\descrizione{Volumina}{Usa questo QR code dal tuo dispositivo mobile per visualizzare o scaricare la partitura completa e le note d’esecuzione}

\begin{pspicture}(1in,1in)
	\psbarcode{https://goo.gl/JCU3Sc}{eclevel=L width=1 height=1}{qrcode}
\end{pspicture}\footnote{https://goo.gl/JCU3Sc}

% \descrizione{Studio Mannaro}{ is a composition for baritone saxophone and live electronics. The saxophone is equipped with a microphone used both for amplification and sound pickup for processing by the computer. The amplified sound together with the processed sounds is then projected in the concert hall using a sound spatialization system based on virtual sources.}

\descrizione{Studio Mannaro}{ per saxofono baritono ed elettronica è stato presentato in prima esecuzione al festival Archipel di Ginevra nel 2012. L’idea poetica del brano è la ricerca dell’unità fra molteplici possibili dissociazioni: il registro estremo grave e estremo acuto del sassofono, le componenti fondamentali e armoniche dello spettro sonoro, la presenza tangibile dello strumento solista e la molteplicità virtuale dei campionamenti e delle sintesi elettroniche. Questo percorso di accettazione dell’identità complessa di un organismo “diverso” e peculiare si realizza attraverso il passaggio fra vari stadi sonori e conseguenti paesaggi immaginari ai quali non sono estranee certe categorie care alla musica di matrice popolare: il drone e il groove.}




\end{multicols}

\clearpage

%----------------------------------------------------------------------------------------

\section*{8 ottobre 2015 -- ore 17:00}

\subsection*{{\small Concerto Audiovisuale} \\
	\textsf{Sala Medaglioni}}

{\fontsize{30}{30} \svolk{\emph{ATTO VII}}}

\subsection*{\textsf{regia del suono di Leonardo Mammozzetti}}

\bigskip

\begin{multicols}{2}

%% !TEX encoding = UTF-8 Unicode
% !TEX TS-program = XeLaTex
% !TEX root = EMU2015_booklet.tex

\acusmatici{Antonio Mazzotti}
{I haven't seen you on the jumbotrons at TimeSquare}{2014}{10'44''}

\acusmatici{Francesc Martí}
{Speech 2}{2015}{7'53''}

\acusmatici{Alfredo Ardia, Sandro L'Abbate}
{Studio N.1}{2015}{2'}

\acusmatici{Carlo Barbagallo}
{Drowning}{2015}{8'00''}

\acusmatici{Roberto Musanti}
{Rotational Chaos}{2014}{3'50''}

\acusmatici{João Pedro Oliveira}
{Hydatos}{2012}{13'00''}

% \descrizione{I haven't seen you on the jumbotrons at TimeSquare}{'I have not seen you on the jumbotrons at TimeSSquare' was realized with the ComputerAidedAlgorithmComposition.It was conceived as a study for the computational models to produce musically meaningful results.The model which I am designing has been adopted as a tool of composition, investigating on the deep connection between sound and emotional meaning.It was implemented in “Mathematica”, “Csound” and Kyma, that uses the Pacarana as audio accelerator. "Processing" for the rendering video.}

\descrizione{I haven't seen you on the jumbotrons at TimeSquare}{Il modello che ho progettato indaga a fondo il legame tra le composizioni multimediali audio/video ed il significato emotivo. Il video e  l'audio in sincronia sono generati utilizzando lo stessa algoritmo. La composizione si sviluppa basandosi su pratiche di tipo algoritmico per progettazione assistita dal computer. L’ insieme di simboli gestiti da una serie di regole inserite nell’algoritmo producono espressioni complesse. Gli errori, imperfezioni e limitazioni dei particolari mezzi compositivi sono elementi costitutivi del pezzo. L’implementazione per il tramite di "Mathematica", "Csound" e’ Kyma’, che utilizza Pacarana come acceleratore audio. "Processing" per il video rendering.}

% \descrizione{Speech 2}{Speech 2 is an experimental audiovisual piece created from a series of old clips from the public affairs interview program The Open Mind. Technically, in this piece, the author has been experimenting how granular sound synthesis techniques, and pseudo-random number generators can be used for audiovisual creative works. All the piece sounds and images come from that series of clip, in other words, no other sound samples or images have been used to create the final result.}

\descrizione{Speech 2 è un brano audio-video sperimentale ricavato da una serie di vecchi video presi da interviste del programma di affari pubblici "The Open Mind". Tecnicamente, in questo brano, l'autore ha sperimentato in quali modi le tecniche di sintesi granularee i generatori di numeri pseudorandomici possono essere utilizzati in lavori creativi audiovisuali. Tutti i frammenti di suono e le immagini provengono da queste serie di video, in altre parole, nessun altro campione di suono o di immagini sono state utilizzate per ricreare il risultato finale.}

% \descrizione{Studio N.1}{A disoriented individual interacts with a virtual space. The body becomes an instrument, therefore able to respond to stimuli according to a specific sound logic. As any computer program, it is a system where the defects, bugs, disruptions and any kind of interferences, lead to a fragmented pace of the narrative itself. This work was born with the intention of creating a connection between what we see and what we hear, it's a study about connections and interactions of these processes. The simple sound material, sine waves and glitches, presents complex internal movements based on beats phenomenon linked to the dynamics of the video.}

\descrizione{Studio N.1}{Un individuo interagisce all’interno di uno spazio virtuale. Il corpo diviene strumento in grado di rispondere agli stimoli, secondo precise logiche sonore. Come ogni programma informatico, esso è un sistema in parte difettoso dove interferenze e interruzioni di vario genere determinano un andatura frammentata della narrazione. Il lavoro è nato con l’intenzione di creare una relazione tra ciò che si osserva e ciò che si ascolta, uno studio sul legame e sull’interazione di questi due aspetti. Il materiale sonoro utilizzato, sinusoidi e glitch, apparentemente semplice e di carattere minimalista, presenta in realtà complessi movimenti interni, dovuti al fenomeno dei battimenti, che seguono le dinamiche del video secondo precisi criteri compositivi.}

% \descrizione{Drowning}{Drowning is a three movements piece written by Carlo Barbagallo for Jean Francois Laporte's Table de Babel. Video directed & performed by Isobel Blank. Drowning: I: Surprise / Involuntary Breath Holding II: Unconsciousness III: Hypoxic Convulsions / Clinical Death}

\descrizione{Drowning}{I: Surprise/Involuntary Breath Holding II: Unconsciousness III: Hypoxic Convulsions/Clinical Death - Drowning, un invito ad annegare, è un brano in tre movimenti composto per la Babel Table, strumento musicale che utilizza l’aria compressa come sorgente di energia, d’invenzione di Jean- Francois Laporte. Ispirata alla fasi dell’annegamento, la composizione del brano è stata commissionata da Productions Totem Contemporain in collaborazione con la Scuola di Musica Elettronica del Conservatorio di Torino. Il video è stato diretto e realizzato dall’artista Isobel Blank. Il brano, nella sua versione dal vivo, è stato già eseguito a Giugno 2015 al Conservatorio G.Verdi di Torino, a Settembre alla Fountain School of Performing Arts della Dalhousie University (Halifax, Canada) e verrà eseguito a Novembre all’università di Huddersfield (UK). Musica: Carlo Barbagallo Video: Isobel Blank Babel Table: Jean-Francois Laporte}


% \descrizione{Rotational Chaos}{"Chaos rotation" is an audiovisual work that explores the relationship between images and sounds. Although the composition is abstract, because it is based primarily on the relationship between the forms, and between them and the sounds, the assembly of the material has the effect of a kind of narrative that we can define Geo/Math fiction.}

\descrizione{Rotational Chaos}{è un lavoro audiovisuale che esplora il rapporto tra suono e immagine. In primo piano, solidi di rotazione, il cui profilo è influenzato dal contenuto armonico di suoni da generatori caotici,  evidenziano il contrasto tra le simmetrie delle loro forme e il timbro sonoro. In secondo piano un sistema particellare funge da "sfondo" alla composizione. Benché si tratti di un lavoro astratto, i materiali grafici scelti e il loro montaggio, anche in rapporto al suono, invitano a una lettura “narrativa”, una sorta di geo/math-fiction.}

% \descrizione{Hydatos}{Hydatos is a greek word that means “water”. This piece is inspired on the first verses of the Old Testament (Genesis Chapter 1:2) “And the Spirit of God moved upon the face of the waters.” The audio part of this piece was commissioned by Gulbenkian Foundation, and was composed in the composer’s personal studio and at the NOVARS Center in Manchester. The video part was done at the composer’s personal studio.}

\descrizione{Hydatos}{Hydatos è una parola greca che significa “acqua”. Questo brano trae ispirazione dai primi versi del Vecchio Testamento (Genesi Capitolo 1:2) “E lo spirito di Dio aleggiava sulla superficie delle acque.” La parte sonora di questo video è stata commissionata dalla Fondazione Gulbenkian, ed è stata composta nello studio personale del compositore e al NOVARS Center a Manchester. La parte video è stata realizzata nello stesso studio personale del compositore.}




\end{multicols}

\clearpage
%----------------------------------------------------------------------------------------

\section*{8 ottobre 2015 -- ore 20:30}

\subsection*{{\small Concerto Acusmatico \& Live Electronics\footnote{ Il concerto verrà trasmesso in diretta streaming da Radio Cemat}} \\
	\textsf{Sala Accademica}}

{\fontsize{30}{30} \svolk{\emph{ATTO VIII}}}

\subsection*{\textsf{Regia del suono di Marco De Martino}}

\bigskip

\begin{multicols}{2}

%CONCERTO 8
Concerto 20.30 (Sala Accademica)

\livel{Enrico Minaglia}
{Physis III - Metal}%7'10''
{per trombone e live electronics}
{trombone}{Raffaele Marsicano}

\descrizione{Physis III - Metal}{This is the third piece of the Physis series, a tribute to Giordano Bruno's philosophy. Physis III  is focused on the "De gli eroici furori dialogue". It aims to depict in music the calm life of the wise man, and its sudden tranformation into the "furioso", a man who struggles to penetrate the misteries of Nature, and in the end will pay this knowledge with his life, like the mythical hunter Actaeon, and Giordano Bruno himself.}

\descrizione{Physis III - Metal}{Si tratta del terzo brano della serie "Physis", un tributo al pensiero di Giordano Bruno. Physis III è basata sul dialogo "De gli eroici furori". Ciò che la musica rappresenta è la vita tranquilla del saggio, e la sua improvvisa trasformazione nel "furioso", un uomo che lotta per penetrare i misteri della natura e che pagherà la sua conoscenza con la sua stessa vita, come la cacciatrice Actaeon, e Giordano bruno stesso.}


\livel{Núria Giménez-Comas}
{No More Words}{2015} %6'27''
{per soprano e live electronics}
{soprano}{Virginia Guidi}

\descrizione{No More Words}{Different concepts and ideas has given some inputs to the global conception of the piece: as information overload and non-useless information or text-speechs. To work with this ideas, I have used an open software of “speech to text” (Mary) of some “historical apologies” (found on the net) with different “artificial voices”. After I have granulated this texts to create voice masses in the space around us, (using a library of Open Music called Prisma). To recreate the idea of different spaces ‘saturation’ I have also introduced to the beginning some fragments of noisy landscapes with spoken voices and resynthesis layers of this soundscapes. This sound waves and masses are combined with a voice line that appears gradually and it’s given by the poem of Edgard Allan Poe. In “A dream within a dream” the poet is dealing in what we perceive as a reality and non-reality.}%aggiungere testo Poe

\descrizione{No More Words}{Idee e concetti differenti hanno contribuito alla concezione globale del brano: il sovraccarico di informazioni, i dati non inutili, il testo parlato. Per lavorare con queste idee, ho usato un software libero "speech to text" (Mary) con alcune "scuse storiche" (trovate in rete), con diverse voci artificiali. Dopo di che ho granulato questi testi per creare masse vocali nello spazio tutt' attorno, (utilizzando una libreria di Open Music chiamata Prisma). Per ricreare l'idea di "saturazione ' dei diversi spazi ho inoltre introdotto all' inizio del brano alcuni frammenti di paesaggi rumorosi con voci recitanti e strati di resintesi di questi paesaggi sonori. Le onde sonore e le masse sono state poi unite ad una linea di voce che appare gradualmente dalla poesia di Edgar Allan Poe. In \textit{Un sogno dentro un sogno}", il poeta si occupa di ciò che percepiamo come realtà e non realtà.}%aggiungere testo Poe


\acusmatico{Silvia Lanzalone}
{eRose}{2013} %9'54''
{acusmatico}

\descrizione{eRose}{The word ‘eRose’ is composed by the words ‘electronic’ and ‘rose and was created for this piece to indicate the contrast between the natural and the virtual. The eRose piece is referred to the woman’s world: women’s bodies are modified and transfigured by digital communication, like ‘eroded roses’. Real sounds are totally transformed by computer to express the feeling of unease, but also the feeling of new discovery. The sounds are repeatedly clipped, ‘eroded’ or, on the contrary, sometimes carefully smoothed and refined.}

\descrizione{eRose}{La parola ‘eRose’ è stata coniata per questo brano attraverso la contrazione delle parole inglesi 'electronic' e 'rose', ed è utilizzata per i suoi eterogenei significati è anche per la struttura del fenomeno esogeno che descrive. Il brano è rivolto ad un universo femminile, di cui viene messo in evidenza il processo di decontestualizzazione che la comunicazione digitale ha intrapreso sull'immagine della donna. La parola 'eRose' Il materiale sonoro è infatti, continuamente eroso, corroso, abraso, oppure a volte, al contrario, anche meticolosamente limato, levigato, patinato. }

\livel{Luciano Azzigotti}
{Vilanos}{2014} %10''
{per due flauti amplificati}
{flauto}{Alessandro Pace, Elena D'Alò}

\descrizione{Vilanos}{When a flower dies, and it’s petals wither, a perfect sphere formes inside. Like the subway, the veins, the umbrellas, like the earthworm’s path... arboreal systems grow regardless of the materials that comprise them, displaying the same geometry. They really come to fill the grooves that already exist in the space. The flute sound is transformed into a swirling turbulence, air that is fragmented and divided into a spiral. Hence the re- feeding of a listening technique transforms holes into resonators, the duo as symmetrical exchange point.}

\descrizione{Vilanos}{Quando un fiore muore, e i suoi petali appassiscono, al suo interno si forma una sfera perfetta. Come la metropolitana, le vene, gli ombrelli, come il percorso sotterraneo di un lombrico... I sistemi boreali crescono senza riguardo verso i materiali che li comprimono, percorrendo la stessa geometria. Arrivano a ricoprire le scanalature che già esistono nello spazio. Il suono del flauto è trasformato in una vorticosa turbolenza, l'aria viene frammentata e divisa in una spirale. Pertanto una tecnica d'ascolto di rialimentazione reciproca trasforma i fori in risuonatori, il duo come un punto di simmetria centrale.}


\livel{Karen Power}
{Deafening silence}{2014} %5'54"
{pianoforte e live electronics}
{pianoforte}{Francesco Ziello}

\descrizione{Deafening silence}{is simply about fusing an acoustic instrument, in this case the piano, with a series of cricket recordings from Australia, Japan, Crete, Laos, Cambodia, Italy and Canada. Field recording has become a large part of my artistic practice and I hope that this is the first work in a series for solo acoustic instrument and natural environments from around the world. All recordings have been made by me while attentively listening to each space.}

\descrizione{Deafening silence}{Il silenzio assordante è semplicemente la fusione di uno strumento acustico, in questo caso il pianoforte, con una serie di registrazioni di grilli dell'Australia, Giappone, Creta, Laos, Cambogia, Italy e Canada. Il Field recording ha preso una larga parte della mia pratica artistica e spero che questo sia il primo lavoro di una serie per solo strumento acustico e ambienti naturali del mondo. Tutte le registrazioni sono fatte da me, mentre ascoltavo attentamente ogni spazio.}


\livel{Pasquale Citera}
{Musica per organi caldi}{2015} %6'
{intonazione e fuga elettroacustica per organo e live elctronics}
{organo}{Giovanni Ubertini}

\descrizione{Musica per organi caldi}{Usa questo QR code dal tuo dispositivo mobile per visualizzare o scaricare la partitura completa e le note d’esecuzione} %QR code partitura - link

\descrizione{Musica per organi caldi}{You can scan this QR code with your mobile device to view or download the complete score and performance notes}%QR code partitura - link





\end{multicols}

\clearpage

%----------------------------------------------------------------------------------------

\section*{9 ottobre 2015 -- ore 17:00}

\subsection*{{\small Concerto Acusmatico in Cupola Ambisonic} \\
	\textsf{Aula Bianchini}}

{\fontsize{30}{30} \svolk{\emph{ATTO IX}}}

\subsection*{\textsf{regia del suono di Luana Lunetta}}

\bigskip

\begin{multicols}{2}

%% !TEX encoding = UTF-8 Unicode
% !TEX TS-program = XeLaTex
% !TEX root = EMU2015_booklet.tex

\acusmatici{Daniel Osorio}
{Spiegelung}{2013}{9'38''}

\acusmatici{Daniel Blinkhorn}
{frostbYte - wildflower}{2014}{13'00''}

\acusmatici{Antonio Carvallo}
{Vri}{2013}{6'59''}

\acusmatici{Kenn Mouritzen}
{Cat-back}{2015}{8'39''}

\acusmatici{Kyle Vanderburg}
{Reverie of Solitude}{2014}{10'00''}

% \descrizione{Spiegelung}{"Spiegelung" (Image on the Mirror) is a work commissioned for the DVD project "Long Sguardo". The work is based on the Stefano AmoresÈs text "Specchiatura" and later became part of the video by visual artist Horkay Istvàn. "Spiegelung", an acousmatic piece for 5.1 surround, was composed from recordings by soprano Franziska Erdmann and narrator Jens Alles, to which was added the digital processing of different sound sources, also recorded for the work. The complexity and length of the text allowed generating different sound perspectives through different forms of sound synthesis, which however, are subordinate to rhythm and sound density of spoken text (German).}

\descrizione{Spiegelung}{"Spiegelung" (immagine a specchio) è un lavoro commissionato per il progetto DVD "Long Sguardo". Il lavoro si basa sul testo di Stefano Amorese "Specchiatura" e, in seguito, è entrato a far parte del video dell'artista visivo Horkay Istvàn. "Spiegelung", un brano acusmatico per surround 5.1, è stato composto con registrazioni del soprano Franziska Erdmann e del narratore Jens Alles, a cui è stata aggiunta l'elaborazione digitale delle diverse sorgenti sonore, anch'esse registrate per il lavoro compositivo. La complessità e la lunghezza del testo hanno permesso la generazione di prospettive sonore differenti attraverso diverse forme di sintesi del suono, che sono comunque subordinate al ritmo e alla densità del suono del testo parlato (in Tedesco).}

%\descrizione{Spiegelung}{„Spiegelung„ es una obra encargada para el proyecto en DVD „Largo Sguardo“ y se basó en el texto de Stefano Amorese („Specchiatura“). Es una pieza compuesta para surround 5.1 a partir de las grabaciones Franziska Erdmann y Jens Alles, a lo que se agregó el procesamiento digital de diferentes fuentes sonoras, también grabadas para la obra. La complejidad y extensión del texto permitió generar diferentes perspectivas sonoras a través de diferentes formas de síntesis de sonido.}

%\descrizione{frostbYte - wildflower}{frostbYte is the last in a cycle of works using field recordings from the high artic region of Svalbard. What was most discernible when recording fragments of glacial ice floating in fjords were the many and varied sonorous ecosystems emanating from underwater, each with its own distinctive personality. In every instance the ice fragments reacted differently to temperature, pressure and other observable phenomena, producing similar, yet unique sonorities. Throughout the work I wanted to capture some of the delicate complexity, as well as the unified symmetries produced through the charismatic, audible ecosystems indelibly linked to each of the naturally formed ice sculptures.} % In order to transcribe, then sculpt these natural carvings into gestures and phrases within the piece I chose to de-construct a number of hydrophone recordings into discrete elements, often organised into families of sound shapes. These typomorphologies were then re-constructed into a variety of gestures, phrases and forms, each of which contained its own attendant ecosystem of sound, much like the original field recordings. From a broader perspective, the resultant phrases are intended to mimic the idea of something that is carved. To my mind the final geometries and patterns sculptured became like those of the short-lived wildflowers that grow in the region, each populating its own unique ecosystem and all subject to the natural forces at play around them.

\descrizione{frostbYte - wildflower}{frostbYte è l'ultimo di un ciclo di opere che utilizzano registrazioni ambientali realizzate nell'alta regione artica di Svalbard. Ciò che è risultato più evidente durante le registrazioni dei frammenti di ghiaccio galleggianti nei fiordi, sono stati i molti e vari ecosistemi sonori sott'acqua, ognuno con la propria distinta personalità. I frammenti di ghiaccio hanno reagito in modo diverso a seconda della temperatura, pressione e altri fenomeni osservabili, producendo sonorità simili, ma uniche.In questo lavoro ho voluto catturare un po ' della delicata complessità e contemporaneamente le simmetrie unificate, prodotte attraverso i carismatici ecosistemi sonori indelebilmente legati a ciascuna delle sculture di ghiaccio formatesi naturalmente.} %Al fine di trascrivere, poi scolpire queste sculture naturali in gesti e frasi all'interno il pezzo, ho scelto di de-costruire una serie di registrazioni con idrofoni in elementi discreti, spesso organizzati in famiglie di forme sonore. Queste tipomorfologie sono state poi ri-costruite in una varietà di gesti, frasi e forme, ognuna delle quali conteneva un proprio ecosistema sonoro, proprio come i field recordings originali.Da una prospettiva più ampia, le frasi risultanti sono destinate ad imitare l'idea di qualcosa che viene scolpito. A mio avviso, le geometrie finali e i modelli scolpiti sono diventati come quelli dei fiori selvatici dalla breve vita che crescono nella regione, i quali popolano il proprio unico ecosistema e sono tutti soggetti alle forze naturali che agiscono intorno ad essi.

% \descrizione{Vri}{The piece is made from the deployment of numerical proportions in charge of controlling the parameters that regulate the synthesis and processing of sound, proportions that somehow guarantee a certain sonority and organization that tends to achieve unity. The piece is made from synthesis subtractive (white noise filtered) and convolution. The resulting sounds are analyzed and their formants with more energy are filtered; finally, an analysis and resynthesis process generates frequency bands that are being shipped, differentially, to the four channels.}

\descrizione{Vri}{Il pezzo è realizzato utilizzando il dispiegamento di proporzioni numeriche incaricate di controllare i parametri che regolano la sintesi e l'elaborazione del suono, proporzioni che in qualche modo garantiscono una certa sonorità e organizzazione, in modo da raggiungere l'unità. Il pezzo è realizzato in sintesi sottrattiva (rumore bianco filtrato) e convoluzione. I suoni risultanti sono analizzati e le formanti con più energia vengono filtrate; infine, un processo di analisi e risintesi genera bande di frequenza che vengono spedite, in modo differenziato, ai quattro canali.}

% \descrizione{Cat-back}{Cat-back is a micro composition on basis of bass clarinet. The recordings were extensively treated in order to give other character while at the same time keeping the powerfull mechanical effect this instrument has.}

\descrizione{Cat-back}{Cat-back è una micro composizione basata su clarinetto basso. Le registrazioni sono state ampiamente trattate in modo da dare un altro carattere al suono, ma al tempo stesso mantenere il potente effetto meccanico dello strumento.}

% \descrizione{Reverie of Solitude}{The piece serves as both an exploration of and a invitation to reverie, providing a space wherein the listener is asked to reconsider their idea of what it means to daydream. At once immersed in a familiar crowd hum, lost among the multitude; it is easy to believe that this daydream is not an expression of solitude, but rather a longing for solitude. And so the piece suggests the pattern of a day dream: the crowd noise giving way to a train, a lazy lawn sprinkler, a contemplative rain storm, a frothing river which becomes a bucolic afternoon on the lake. Each vignette is a self-contained narrative wherein to consider solitude in a natural context. The metaphor of water and the alternating themes of movement and respite invite the listener to reflect on the purpose of a daydream: to escape, to pacify, or to enrich a perfect moment. After having their attention turned to the daydream they themselves have been lulled into, the listener is returned to the crowd hum having established a personal sense of solitude within the piece and within the audience. Program Note by Walter Jordan}

\descrizione{Reverie of Solitude}{Il brano può essere considerato sia come un'esplorazione che un invito al fantasticare, fornendo uno spazio in cui l'ascoltatore è invitato a riconsiderare l' idea di ciò che significa sognare ad occhi aperti. All'inizio l'ascoltatore è immerso in un familiare ambiente con ronzio di folla, perso tra la moltitudine; è facile credere che questo sogno ad occhi aperti non sia espressione della solitudine, ma piuttosto un desiderio di solitudine. E così la composizione suggerisce il modello di una fantasticheria: il rumore della folla cede il posto a quello di un treno, un pigro sistema d'irrigazione per prati, una contemplativa tempesta di pioggia, la schiuma di un fiume che diventa un pomeriggio bucolico sul lago. Ogni vignetta è una narrazione autonoma che prende in considerazione la solitudine in un contesto naturale. La metafora dell'acqua e l'alternarsi di  tematiche di movimento e di tregua invitano l'ascoltatore a riflettere sullo scopo del sogno: la fuga, per pacificare, o per arricchire un momento perfetto. Dopo essere stato cullato dal sogno ad occhi aperti, l'ascoltatore viene restituito al ronzio della folla dopo aver stabilito un personale senso di solitudine all'interno del brano e nel pubblico. Note di Walter Jordan}


\end{multicols}

\clearpage
%----------------------------------------------------------------------------------------

\section*{9 ottobre 2015 -- ore 20:30}

\subsection*{{\small Concerto Acusmatico \& Live Electronics\footnote{ Il concerto verrà trasmesso in diretta streaming da Radio Cemat}} \\
	\textsf{Sala Accademica}}

{\fontsize{30}{30} \svolk{\emph{ATTO X}}}

\subsection*{\textsf{Regia del suono di Francesco Ziello}}

\bigskip

\begin{multicols}{2}

%% !TEX encoding = UTF-8 Unicode
% !TEX TS-program = XeLaTex
% !TEX root = EMU2015_booklet.tex

\livel{Javier Alejandro Garavaglia}
{Hoquetus}{9'}
{for Soprano saxophone and multi-track electronics}
{2005-6}

\brano{Gustavo Delgado}
{Permanente e transitorio}{7'20''}
{acusmatico}{2015}\\

\livel{Sylvano Bussotti}
{Ultima RARA}{8'}
{per un chitarrista recitante}
{versione 2015}

\livel{James Dashow}
{Soundings in Pure Duration n.7}{10'}
{per sassofono contralto e live electronics}
{2015}

\livel{Maurizio Pisati}
{SENTI!}{15'}
{chitarra, percussione, orchestra d’archi e live electronics}
{versione 2015}

%\vspace{5mm}

\esecutore{sassofoni}{Enzo Filippetti}
\esecutore{chitarra}{Arturo Tallini}
\esecutore{percussioni}{Gianluca Ruggeri}
\esecutore{orchestra d'archi}{Orchestra Europa}

% \descrizione{Hoquetus}{The first version of the piece was commissioned by Prof. Esther Lamneck (NYU), who also premiered it at the 14th Florida Electroacoustic Music Festival on the Tárogató (April 2005). The current version is a further development, with changes in the instrumental part (including its adaptation for Saxophone) as well as in the electronics, which have been completely reprogrammed since the world premiere.Basically, the work tries to recreate in a different context the medieval technique of the Hocket. Being this technique polyphonic, the Tárogató part delivers a sort of “hiccup” technique with itself by the use of different registers. Moreover, the computer makes a Hocket-like response to the Tárogató at certain moments.}

\vspace{-5mm}
\descrizione{Hoquetus}{La prima versione di questo brano fu commissionata dal prof. Esther Lamneck (NYU), e fu premiato alla quattordicesima edizione del Florida Electroacoustic Music Festival (Aprile 2015). La versione revisionata ha avuto diversi sviluppi, con il cambio da parte degli strumenti (compreso il suo adattamento per sassofono) non che nell'elettronica, dove è stata completamente riprogrammata rispetto alla prima mondiale. Il lavoro ha cercato di ricreare in un differente contesto la tecnica medievale dell' Hoquetus. Essendo una tecnica polifonica, le parti di Tárogató fornisce una sorta di tecnica del "singhiozzo" mediante l uso di diversi registri. Inoltre, in certi momenti  il computer esegue una risposta simil-hoquetus al Tárogató.}

% \descrizione{Permanente e transitorio}{The composition deals with two seemingly opposite concepts in a such closer relationship: permanent and temporary. By the combination of different mixing and editing techniques very often used in the cinema and video games sfx, there were created very complex "impact sounds" classified by their dynamic structure as short (temporaries). Resonance frequencies have been isolated and emphasized from those materials in order to obtain stable sounds (permanents) to be used as starting points or bridges to the original impact sounds and viceversa.}

\descrizione{Permanente e transitorio}{Mediante il montaggio e la combinazione di diverse tecniche di missaggio e sound design utilizzate al cinema, ai videogiochi e alla musica elettronica, sono stati creati numerosi materiali di tipo “impact sounds” (materiale di tipo transitorio) da cui sono estratte delle risonanze (materiale di tipo permanente) dinamicamente modulate in frequenza (FM) su una patch in MaxMSP creata dall’autore.}

% \descrizione{Ultima Rara}{eng}

\descrizione{Ultima Rara}{Ultima Rara è pensata “per voce e da una a tre chitarre” la versione di questa sera riunisce le 4 figure e infatti viene proposta come versione “per un chitarrista recitante” così come proposta da Arturo Tallini con l’approvazione del compositore. Il testo e conseguentemente la musica hanno un carattere multiforme, fortemente teso all’espressione intima, seppure non disgiunta da un’ironia che qua e là fa capolino. Il brano ci racconta la fine di una storia passando dal dolore, al delirio alla malinconia. Il tipo di recitazione non è indicata in alcun modo dunque la versione di stasera è stata interamente elaborata dall’interprete.}

% \descrizione{Soundings in Pure Duration n.7}{The seventh in the series of Soundings in Pure Duration provides a wealth of timbres and trajectories in octophonic space designed to place the alto saxophone soloist in the middle of varying degrees of transformative motion, musical and  physical.  No, the saxophonist doesn't fly around the hall, but the three-dimensional temporally structured spatialization of the constantly evolving sounds creates the impression of everything, including the soloist, moving in contrappuntal complexities.}

\descrizione{Soundings in Pure Duration n.7}{Il settimo brano della serie Soundings in Pure Duration propone un'abbondanza di timbri e tragitti nello spazio ottofonico, creati per collocare il sassofono contralto al centro di una varietà di movimenti trasformativi, musicali e fisici.  No, il sassofonista non vola attraverso la sala, ma la spazializzazione dei suoni in evoluzione è strutturata  temporalmente nelle tre dimensioni così da creare l'impressione che tutto, incluso il solista, si muova nelle complessità contrappuntistiche.}

% \descrizione{SENTI!}{In SENTI! (listen up!) the sounds and gestures born in Senti? (can you hear?) and Sì. Cos’è? (Yes. What is this?) came to an end, closing a triptyque made in the two previous operas. Now the final disposition of the compositions is transformed into an essential dialogue: «Senti? Sì, cos’è? SENTI!». The triptyque is prompted by a short story by August Strindberg “Stora grusharpan” here adapted at the composer will. This leads us into the sea bottom, and from here it tells of sounds never heard before: a Piano falls into water, a storm deforms it and the fishes explore it lightly touching it’s string, while on the shore two guys are listening. The title of the first track –Senti?– it’s the question that one of the two asks to the other after hearing that sounds. In the second track the friend answers, asking too “Sì. Cos’è?”.  In the latter part, after wondering briefly on the nature of the sounds, no one asks anything or waits for an answer, until one of the two states: “sì, ma senti!”. From the brief dialogue of the two guys, still able to appreciate the marvel, is born the third score, which owes his shape to the imaginary breathe of things and sounds sinked in unreachable places.}

\descrizione{SENTI!}{In SENTI! giungono alla conclusione i gesti e alcune sonorità nate in Senti? e Sì. Cos’è?, i due primi pezzi di un trittico. Ora quindi l’ordine finale delle tre composizioni è diventato anche un dialogo essenziale: «Senti? Sì, cos’è? Senti!» Il trittico prende le mosse dalla novella di August Strindberg “Stora grusharpan” qui liberamente adattata. Ci conduce su un fondale marino e da lì racconta di sonorità mai sentite: un Pianoforte cade in acqua, una tempesta lo deforma e i pesci lo esplorano sfiorando le corde, mentre sulla riva due ragazzi ascoltano. Il titolo del primo brano –Senti?– è la domanda che uno rivolge all’altro dopo aver udito quei suoni. Nel secondo pezzo l’amico risponde, a sua volta chiedendo: “Sì, cos’è?”. Infine qui, nel terzo brano, dopo un breve scambio di supposizioni sulla natura di quei suoni, nessuno più domanda o attende risposta, ma uno dei due formula la frase forse più importante, l’esortazione ad ascoltare: “sì, ma senti!”. Dal breve dialogo dei due ragazzi che ancora sanno godere della meraviglia, nasce questa terza partitura, che deve la sua forma al respiro immaginario delle cose e dei suoni sprofondati in un luogo di solito inaccessibile.}


\end{multicols}

\clearpage
%----------------------------------------------------------------------------------------

\section*{12 ottobre 2015 -- ore 18:00}

\subsection*{{\small Concerto Acusmatico \& Live Electronics\footnote{ Il concerto verrà trasmesso in diretta streaming da Radio Cemat}} \\
	\textsf{Auditorium Ennio Morricone}}

{\fontsize{30}{30} \svolk{\emph{ATTO XI}}}

\subsection*{\textsf{regia del suono di Massimiliano Mascaro}}

\bigskip

\begin{multicols}{2}

%% !TEX encoding = UTF-8 Unicode
% !TEX TS-program = XeLaTex
% !TEX root = EMU2015_booklet.tex

\livel{Frei}
{Poetics}{16'50''}
{per due laptop}
{2015}

\livel{Christian Banasik}
{Ik ́}{9'}
{per flauto e live electronics}
{2013}

\livel{Marco Marinoni}
{Finita è la terra}{4'16''} %
{pianoforte e live electronics}
{2015}

\livel{Domenico De Simone}
{A[LIVE]}{5'30''}
{pianoforte, vibrafono  ed elettronica su supporto}
{2015}

\brano{Benjamin D. Whiting}
{Illumina! Arabidopsis thaliana}{9'18''}
{acusmatico}{2015}\\

\esecutore{flauto}{Elena D'Alò}
\esecutore{pianoforte}{Sara Ferrandino}
\esecutore{vibrafono}{Matteo Rossi}

\noindent \textsc{Frei}:\\
\esecutore{due laptop}{Francesco Bianco, Paolo Gatti}

% \descrizione{Poetics}{Work based on the concepts of entropy, redundancy, and noise and how these play a role on messages and information in a musical discourse. From an aesthetic point of view, everything is resolved, here, in an alternation between punctiform moments, almost static, and growing dynamic more or less sudden.}

\descrizione{Poetics}{Lavoro basato sull'improvvisazione giocata sull'alternanza fra momenti puntiformi, quasi statici, e dinamiche di crescendo più o meno improvvise. Altro tema trattato è il rapporto fra comunicabilità e incomunicabilità tramite l'indagare delle modalità di trasmissione dei messaggi musicali e vocali e i concetti di entropia, ridondanza e di rumore.}

% \descrizione{Ik ́}{\textit{Ik ́} is the name of the 2nd day in the ritual calendar of the Maya associated with breath and wind.The liturgical year consisted of twenty cycles and their glyphs, each of them thirteen days long, and had 260 days in all. For the tone material I use a contemporary folk song from Central America and a virtually generated original song of the Maya.The flute score consists of virtuoso abstracted variations mixed with the graphical-audio analysis of this particular glyph. It is divided in 20 parts.}

\descrizione{Ik ́}{\textit{Ik ́} è il nome del secondo giorno del calendario rituale dei Maya, associato al respiro e al vento. L'anno liturgico era composto da venti cicli e i loro glifi, ognuno dei quali durava tredici giorni, avendo in totale 260 giorni. Per il materiale tonale ho usato una canzone popolare contemporanea dell'America Centrale, e una canzone dei Maya generata virtualmente. La parte del flauto consiste in variazioni astratte miste all'analisi grafica-audio di questo particolare glifo. È diviso in 20 parti.}

% \descrizione{Finita è la terra}{is a small formal experiment that makes use of three piano sound objects, variants of a single musical image, which remain identical to themselves during the entire piece, arranged in a formal course partly iterative and partly structured through microvariations. These objects interact along the time axis with material fixed on support, agglomerations deriving from their possible electroacoustic processing – a live-electronics frozen, or even more than a live-electronics, an image of it, captured and distilled through an alchemical process which aims to turn something that was only movement, dispersion, probability (grain clouds applied to the three sound events of origin) in its \greco{εἴδωλον}, playing this way with the perception and memory in a continuum consisting of appearance / disappearance / memory / expectation / desire.}

\descrizione{Finita è la terra}{è un piccolo esperimento formale che fa uso di tre oggetti sonori pianistici, varianti di un'unica immagine musicale, che permangono identici a se stessi durante tutto il brano, disposti secondo un decorso formale in parte iterativo e in parte microvariato. Tali oggetti interagiscono lungo l’asse del tempo con materiali fissati su supporto, agglomerazioni derivanti da una loro possibile processazione elettroacustica – un live-electronics ghiacciato, o ancora, più che un live-electronics, una sua immagine, catturata e distillata, mediante un processo alchemico volto a trasformare qualcosa che era solo movimento, dispersione, probabilità (le nubi di granulazioni applicate ai tre eventi sonori di partenza) nel suo \greco{εἴδωλον}, giocando in questo modo con la percezione e la memoria in un continuum fatto di apparizione/sparizione/ricordo/attesa / desiderio.}

% \descrizione{A[LIVE]}{I AM NOT THERE \\ Do not stand at my grave and weep. \\ I am not there. I do not sleep. \\ I am a thousand winds that blow. \\ I am the diamond glints on snow. \\ I am the sunlight on ripened grain.  \\ I am the gentle autumn rain. \\ When you awaken in the morning’s hush \\ I am the swift uplifting rush \\ Of quiet birds in circled flight. \\ I am the soft stars that shine at night. \\ Do not stand at my grave and cry; \\ I am not there. I did not die. \\ Mary Elizabeth Frye}

\descrizione{A[LIVE]}{IO NON SONO LÌ \\ Non piangere sulla mia tomba.  \\ Io non sono lì. Io non sto dormendo.  \\ Io sono mille venti che soffiano. \\ Sono lo scintillio del diamante sulla neve. \\ Sono il sole che brilla sul grano maturo. \\ Sono la lieve pioggia d'autunno.  \\ Quando tu ti svegli nel silenzio del mattino, \\ io sono il rapido fruscio \\  degli uccelli che volano in cerchio. \\ Io sono la piccola stella che brilla di notte. \\ Non piangere sulla mia tomba;  \\ io non sono lì, la mia anima non è morta.   \\ (libera traduzione della poesia I AM NOT THERE di Mary Elizabeth Frye)}

% \descrizione{Illumina! Arabidopsis thaliana}{This piece represents the ongoing artistic and scientific collaboration between genomic biologist Aleel K. Grennan and myself. Grennan is studying the rate of photosynthesis between a natural wild type of Arabidopsis thaliana leaf and three genetically engineered mutants with different sizes of chloroplasts. I designed the majority of the sonic material in DISSCO and KYMA, incorporating Grennan’s data into several parameters, thus creating a wealth of stylized sounds.}

\descrizione{Illumina! Arabidopsis thaliana}{Questo brano rappresenta l'attiva collaborazione artistica e scientifica tra il biologo genetico Aleel K. Grennan. Grennan sta studiando l'indice di fotosintesi tra un tipo naturale selvatico di foglia Arabidopsis thaliana e tre diverse mutazioni geneticamente modificate con diverse grandezze di cloroplasti. La maggior parte dei materiali sonori sono stati progettati con DISSCO and KYMA, incorporando i dati di Grennans attraverso alcuni parametri, sono stati creati così un'abbondanza di suoni stilizzati.}




\end{multicols}

\clearpage
%----------------------------------------------------------------------------------------

\section*{13 ottobre 2015 -- ore 18:00}

\subsection*{{\small Concerto Acusmatico \& Live Electronics\footnote{ Il concerto verrà trasmesso in diretta streaming da Radio Cemat}} \\
	\textsf{Auditorium Ennio Morricone}}

{\fontsize{30}{30} \svolk{\emph{ATTO XII}}}

\subsection*{\textsf{Regia del suono di Francesco Bianco}}

\bigskip

\begin{multicols}{2}

%% !TEX encoding = UTF-8 Unicode
% !TEX TS-program = XeLaTex
% !TEX root = EMU2015_booklet.tex

\livel{Jorge García del Valle Méndez}
{no sun, no moon}{9'20''}
{per flauto basso ed elettronica su supporto}
{2012}

\livel{Giovanni Costantini}
{Traccia sospesa}{7'35''} 
{pianoforte ed elettronica su supporto}
{2015}


\brano{James Andean}
{Déchirure}{9'58''}
{acusmatico}{2013}\\

\livel{Mario Mary}
{Rock}{8'30''}
{per pianoforte ed elettronica su supporto}
{2013}

\livel{Inhorep}
{Inhorep@emufest}{15'}
{laptop e autocostruiti}
{2015}
\\

\esecutore{flauto basso}{Alessandro Pirchio}
\esecutore{pianoforte}{Sara Ferrandino}

\noindent \textsc{Inhorep}:\\
\esecutore{chitarra preparata e laptop}{Davide Palmentiero}
\esecutore{strumenti autocostruiti e laptop}{Giuseppe Pisano}
\esecutore{laptop}{Massimo Varchione}

%\vspace{12mm}

% \descrizione{no sun, no moon}{no sun, no moon immerses into a world without reference points, where the reality is considered by two different points of view, and we do not know which is the real and which is not. In no sun, no moon, the bass flute and the electronics are the two worlds, reality and parallel reality. Both are the same thing and simultaneously its opposite, interacting and reacting one another. The raw materials for the electronics are exclusively bass flute samples. Honorary Mention at the CICEM 2014.}

\descrizione{no sun, no moon}{Con \textit{no sun, no moon} siamo immersi in un mondo senza punti di riferimento, in cui la realtà è considerata da due diversi punti di vista, e non sappiamo quale sia la reale e quale non lo sia. Il flauto basso e l'elettronica sono i due  mondi; realtà e realtà parallela. Entrambi sono la stessa cosa e contemporaneamente il loro contrario, interagiscono e reagiscono tra di loro. Le materie prime dell'elettronica sono esclusivamente campioni di flauto basso. Mensione d'onore al Cicem 2014.}

%\descrizione{Traccia sospesa}{The piece evokes events, places, sounds and feelings related to the First World War. The piano is played in a "non-classical" way, exploring new sonorities useful to convey in the listener pain, dismay, disbelief. The electronics consists in a textures created through elaborations of piano sounds: an alter ego with which the piano can dialogue, in an atmosphere suspended between memories and uncertainties.}

\descrizione{Traccia sospesa}{Il brano vuole evocare avvenimenti, luoghi, suoni e sentimenti legati alla prima guerra mondiale. Il pianoforte diventa strumento utile a trasmettere dolore, sgomento, incredulità, mediante un utilizzo “non classico” che ne esplora sonorità nuove. La parte elettronica è costituita da una trama di tracce sonore realizzate al computer attraverso elaborazioni di suoni di pianoforte: un alter ego con cui dialogare, ricercando e ricordando, in un’atmosfera sospesa e di continua incertezza.}

% \descrizione{Déchirure}{Déchirure: a tearing, a painful separation... This piece involves a number of 'déchirures', both musical as well as figurative, although the only literal 'tearing' is saved for the final phrase. This work was composed for Presque Rien 2013, in which it received a Special Mention. For this project, sounds from Luc Ferrari's archives were made available to composers for the composition of new works; all sound materials used in the piece are sourced and developed beginning from these recordings.}

\descrizione{Déchirure}{Déchirure: uno strappo, una dolorosa separazione…Questo brano comporta una serie di "déchirure" (strappi), sia musicali, così come figurativi, anche se l'unico letterale "strappo" viene lasciato per la frase finale. Questo lavoro è stato composto per Presque Rien 2013, in cui ha ricevuto una Menzione Speciale. Questo progetto, i suoni sono stati resi disponibili da archivi di Luc Ferrari ai compositori per la composizione di nuove opere; tutti i materiali sonori utilizzati nel brano sonori sono di provenienza e sviluppati a partire da queste registrazioni.}

% \descrizione{Rock}{It is a kind of homage to progressive rock, also called symphonic rock, which appeared in the 70s, and continues to accompany me in my inner world, though my compositions are always distinctly contemporary aesthetics. In my teens, groups like Yes, Pink Floyd, Emerson Lake & Palmer, injected me the virus of electronic sound, before than I discovered the existence of contemporary and electroacoustic music. "Rock" is not a rock, but is inspired by the energy and character of the music of my youth. Moreover, this work is based on a virtuous dialogue between the instrumentalist and the electroacoustic part.}

\descrizione{Rock}{Si tratta di una sorta di omaggio al rock progressivo, chiamato anche sinfonico, che è apparso negli anni '70, e continua ad accompagnarmi nel mio mondo interiore, anche se le mie composizioni sono sempre un'estetica decisamente contemporanea. Nei miei adolescenza, gruppi come Yes, Pink Floyd, Emerson Lake \& Palmer, ho iniettato il virus del suono elettronico, prima di scoprire l'esistenza di musica contemporanea ed elettroacustica.“Rock” non è una roccia, ma si ispira l'energia e il carattere.}

% \descrizione{Inhorep@emufest}{\textit{INHOREP} is a project of radical electroacoustic improvisation. The trio (first active as Improvviso) was formed in the course of Electronic Music of the  Conservatory of Naples under the impulse of  Elio Martusciello. The members immediately focused themselves on live exhibition, participating in several Italian festivals. Their performance is based on the interplay between the three musicians (guests are always welcome), strengthened through rehearsals and on the the attitude of listening to each other. The musicians improve theirs instruments with a DIY (Do it Yourself ) philosophy, recovering instruments and other sound objects, using the circuit bending and other techniques.}

\descrizione{Inhorep@emufest}{\textit{INHOREP} è un progetto di improvvisazione radicale elettroacustica. Il trio (prima attivo come Improvviso) si è formato nella classe di Musica Elettronica del Conservatorio di Napoli sotto l'impulso del M° Elio Martusciello. Fin da subito il gruppo ha deciso di puntare sull'esibizione dal vivo, partecipando a diversi festival italiani. Le loro performance si basano sull'interazione tra i tre musicisti (e gli eventuali ospiti, sempre assai graditi), affinata attraverso le prove e l'attitudine all'ascolto reciproco. Individualmente, gli strumentisti costruiscono e migliorano costantemente i loro strumenti attraverso la filosofia D.I.Y. (Do it Yourself), recuperando strumenti e altri oggetti sonori, ricorrendo alle possibilità del circuit bending ed altre tecniche.}



\end{multicols}

\clearpage




%----------------------------------------------------------------------------------------

% ---------------------------------------------------------------------------------------------
\section*{ }

\subsection*{\textsf{Autori ed Esecutori}\\}

{\fontsize{30}{30} \svolk{\emph{Biografie}}}

\bigskip

\begin{multicols}{2}

\biografia{Natasha Barrett}{Compositrice freelance che lavora con la musica, la ricerca e l'uso creativo del suono. Ha conseguito in Inghilterra master e dottorato in composizione elettroacustica, dopo di che, nel 1999, si trasferisce in Norvegia, dove ha vissuto. Si concentrò sulla musica acusmatica e sulla musica strumentale con live electronics. Dal 1999 il suo lavoro con il suono si è ampliato fino a comprendere sound-art, installazioni sonore-architettoniche, tecniche interattive, la collaborazione con i progettisti e scienziati sperimentali e performance live e l'improvvisazione. Esempi recenti di questo includono l'uso di dati scientifici e dei processi geologici in sound-art, composizione spaziale per altoparlante a matrice emisferica e uno speciale interesse nell’HOA, e il suo terzo progetto d’installazione con il gruppo Oceano Design Research Association.

Le sue opere sono eseguite e commissionate in tutto il mondo e ha ricevuto numerosi riconoscimenti , in particolare il Music Prize Consiglio nordico (Norden / Scandinavia, 2006), Giga - Hertz Award (Germania, 2008), Edvard Prize (2004, Norvegia), Noroit- Leonce Petitot (Arras, Francia, 2002 e 1998), Bourges International Electroacoustic Music Awards (Francia del 2001, 1998 e 1995), Musica Nova (2001), IV CIMESP 2001 Concours Scrime, (Francia 2000), International Electroacoustic Creation Competition of Ciberart (Italia 2000), Concours Luigi Russolo (Italia 1995 e 1998), Prix Ars Electronica (Linz , Austria 1998), 9th International Rostrum for electoacoustic music (2002). Le sue installazioni includono un lavoro importante per la Commissione dello Stato Norvegese per l'arte negli spazi pubblici}


\end{multicols}

\clearpage

\section*{ }

%\subsection*{\textsf{}\\}

\hyphenation{Michele Andreotti Guido Capotosto Federico Coderoni Gianmarco Costa
Marco De Martino Simone Giudice Matteo Ilardo Leonardo Mammozzetti
Danilo Marro Massimiliano Mascaro Alessandro Pacetta
Federico Paganelli Ivo Papadopoulos Ivano Pecorini Susanna Rimondotto Federico Ripanti}

\begin{center}

{\fontsize{30}{30} \svolk{\emph{Organizzazione}}}
\medskip

{\fontsize{12}{12} \textsf{Dipartimento di Musica Elettronica}}\\
\vspace{.5cm}

\textbf{\textit{Comitato artistico EMUfest}}\\
Nicola Bernardini, Michelangelo Lupone, Alfredo Santoloci, Franco Sbacco
\medskip

\textbf{\textit{Comitato organizzatore}}\\
Francesco Bianco, Elena D'Alò, Paolo Gatti, Marco Giordano, Virginia Guidi, Luana Lunetta, Massimo Massimi, 	Luigi Pizzaleo, Federico Scalas, Giuseppe Silvi, Anna Terzaroli, Francesco Ziello
\medskip

\textbf{\textit{Supervisione tecnica}}\\
Federico Scalas
\medskip

\textbf{\textit{Responsabili di palco}}\\
Luana Lunetta, Massimo Massimi
\medskip

\textbf{\textit{Regia del suono}}\\
Giuseppe Silvi
\medskip

\textbf{\textit{Regia Conferenze}}\\
Anna Terzaroli
\medskip

\textbf{\textit{Tecnico di registrazione}}\\
Federico Coderoni
\medskip

\textbf{\textit{Tecnici luci}}\\
Massimiliano Mascaro, Simone Giudice
\medskip

\textbf{\textit{Ufficio Stampa}}\\
Francesco Bianco, Paolo Gatti
\medskip

\textbf{\textit{Staff esteso}}\\
Michele Andreotti, Guido Capotosto, Federico Coderoni, Gianmarco Costa,
Marco~De~Martino, Simone Giudice, Matteo Ilardo, Leonardo Mammozzetti,
Danilo Marro, Massimiliano Mascaro, Alessandro Pacetta,
Federico Paganelli, Ivo Papadopoulos, Ivano~Pecorini, Susanna Rimondotto, Federico Ripanti
\end{center}


%\begin{figure}[!h]
%\centering
%\includegraphics[width=6.4cm]{loghi.jpg}
%%\caption{}
%%\label{fig2_42}
%\end{figure}

\end{document}
